\chapter{Normal Labor and Vaginal Delivery}

\section*{Overview}
Normal labor and vaginal delivery is the process by which uterine contractions open the cervix and the baby is delivered through the vagina, followed by delivery of the placenta.

\section*{How Labor Progresses}
Labor is commonly described in three stages: cervical dilation and effacement, birth of the baby, and delivery of the placenta. The mother and baby are monitored throughout labor, and pain relief may include breathing techniques, medicines, or neuraxial anesthesia.

\section*{Preparation and Care}
The care team reviews the pregnancy history, checks the baby's position and heart rate, and monitors maternal vital signs and contractions. A birth plan may address support persons, mobility, pain relief, and newborn care, while remaining flexible for safety.

\section*{Risks and Complications}
Possible complications include prolonged labor, heavy bleeding, infection, tears of the birth canal, shoulder dystocia, and changes in the baby's heart rate. Forceps, vacuum assistance, or cesarean delivery may be needed if spontaneous delivery is not safe.

\section*{Recovery}
After delivery, the mother is monitored for bleeding, blood pressure, pain, and urination. The newborn receives an examination and routine preventive care. Recovery and follow-up depend on the delivery and any complications.

\section*{References}
\begin{enumerate}
  \item MedlinePlus. Labor and delivery. U.S. National Library of Medicine; 2025.
  \item Current Medical Diagnosis and Treatment. McGraw-Hill; 2025.
\end{enumerate}
\clearpage
